\addcontentsline{toc}{chapter}{Resumo}

\begin{center}
    \huge{{\bf Resumo}}
    \vspace{2cm}

\end{center}

\begin{sloppypar}

    Este trabalho apresenta um \textit{pipeline} robusto para detecção automática de estados operacionais de equipamentos industriais, utilizando algoritmos de aprendizado não supervisionado baseados em K-Means. O sistema processa dados de sensores IoT de equipamentos elétricos e mecânicos, identificando automaticamente estados LIGADO/DESLIGADO com 100\% de acurácia. A metodologia implementa um \textit{pipeline} completo que inclui coleta de dados, processamento temporal, interpolação adaptativa, normalização e classificação automática via clusterização. Uma contribuição significativa é a demonstração prática de que soluções simples são superiores a abordagens complexas baseadas em redes neurais para problemas de classificação binária industrial. O sistema foi validado com dados reais de cinco equipamentos industriais, processando mais de 5 milhões de registros. Os resultados demonstram robustez, interpretabilidade e eficiência.

    \vspace{1cm}

    \noindent \textbf{Palavras-chave:} Detecção de Estado Operacional, Equipamentos Industriais, IoT Industrial, K-Means, Aprendizado Não Supervisionado, Classificação Automática, Análise Temporal, Manutenção Preditiva, Indústria 4.0.
\end{sloppypar}

